\bookchapter{第三十章：スクリーンミラーリング、プレゼンとカーナビ連携}

\chapterintro{スマホ画面は大きくても7インチだが、テレビは55インチ、プロジェクターは100インチだ。スマホ画面を大画面に“投げる”と、体験は全く次元が異なる。この章ではスクリーンキャスト、プレゼン、カースル連携の三つの高頻度シーンを明確にし、スマホの能力をより大きな画面に拡張する方法を説明する。}

\booksection{一、ワイヤレスキャストの三大陣営——混同しないように}

ワイヤレスキャストはどれも“スマホ画面をテレビに”のように見えるが、基盤の技術は三種類あり、互いに互換性がない：

\begin{center}
\begin{tabularx}{\linewidth}{|X|X|X|}
\hline
\textbf{技術} & \textbf{使用者} & \textbf{特徴} \tabularnewline
\hline
AirPlay & iPhone/iPad/Mac & Appleエコシステム、体験最も滑らか。Apple TVまたはAirPlay対応テレビが必要 \tabularnewline
Miracast & Androidスマホ & Wi-Fi Direct直接接続、ルーター不要。Androidスマホの"ワイヤレスキャスト"はデフォルトでこれを使用 \tabularnewline
DLNA & 全プラットフォーム & メディアファイル（動画/音楽/画像）のみプッシュ、画面ミラーリングなし。テレビに動画を送って再生するのに適している \tabularnewline
\hline
\end{tabularx}
\end{center}

\booksubsection{各プラットフォームの操作方法：}

\begin{itemize}
  \item \textbf{iPhone}：コントロールセンター > 画面ミラーリング > テレビ/Apple TVを選択。テレビがAirPlayに対応していない場合、Apple TVボックスまたはAirPlay対応のテレビスティックを購入できる。
  \item \textbf{Android}：通知パネルをプルダウン > "ワイヤレスキャスト"/"スマートキャスト"/"マルチスクリーンインタラクション"ボタンを見つける > テレビを選択。スマホとテレビが同じWi-Fiに接続している必要がある。
  \item \textbf{汎用ソリューション}：テレビが古すぎてワイヤレスキャストに対応しない場合、ChromecastやXiaomiボックス/Huaweiボックスを買ってテレビのHDMI端子に挿せばキャスト受信ができる。
\end{itemize}

\booksection{二、有線キャスト——より安定、ゼロ遅延}

ワイヤレスキャストは便利だが遅延があり、ゲームやプレゼンでは体験が良くない。有線キャストはこの問題を解決できる：

\begin{itemize}
  \item \textbf{iPhone}：Lightning/USB-C to HDMI変換ケーブルが必要（約100-200元）。接続するとスマホ画面がテレビ/プロジェクターのHDMI端子に直接出力される。
  \item \textbf{Android}：一部のスマホはUSB-C to HDMI直接接続に対応（スマホがDP Alt Modeをサポートする必要あり、大部分のフラッグシップが対応）。USB-C to HDMIケーブルを一本買えばよい。
  \item \textbf{会議室シーン}：頻繁にスマホでプレゼンする場合は、バッグにHDMI変換ケーブルを常備することを推奨。会議室のテレビ/プロジェクターはほぼHDMI端子があり、挿せば使える。ワイヤレスキャストをあれこれするよりはるかに頼りになる。
\end{itemize}

\booksection{三、スマホがパソコンに——デスクトップモード}

一部のAndroidスマホは“デスクトップモード”に対応——モニターに接続すると、パソコンのデスクトップに似たUIを出力し、マウスとキーボードで操作できる：

\begin{itemize}
  \item \textbf{Samsung DeX}：Samsungフラッグシップの看板機能。モニター接続後にデスクトップUIが起動し、マルチウィンドウ操作が可能。本当のパソコンに最も近い体験。Bluetoothキーボード/マウスと組み合わせれば、出張時にノートPCを持たずに済む。
  \item \textbf{Huaweiデスクトップモード}：Huawei Mateシリーズ/Puraシリーズがモニター接続後に類似機能あり。マルチウィンドウ、ファイル管理、オフィスAppに対応。
  \item \textbf{Motorola Ready For}：Motoフラッグシップも対応、機能はDeXに類似。
  \item \textbf{：ハイエンドプロフェッショナル制作。iPhone 15 ProシリーズはApple ProRes動画撮影に対応（最大4K 60fps）しますが、ProRes RAW形式は現状外部録画デバイスとの組み合わせが必要で、スマホ単体での直接撮影は未対応。}：ライトなオフィス作業（メール返信、文書編集）、PPTプレゼン、大画面で動画視聴。ヘビーなExcel操作や専門的なデザインには不適。
\end{itemize}

\booksection{四、カースル連携——運転中のスマホの正しい使い方}

運転中のスマホ操作は違法だが、ナビ、音楽、電話は必須だ。カースル連携システムで車載画面上でスマホのコア機能を使い、スマホは置いて触らないでおける：

\begin{itemize}
  \item \textbf{Apple CarPlay}：iPhoneユーザーの第一選択。USBまたはワイヤレスでカースルに接続、車載画面にiPhoneのシンプルUIが表示される——ナビ（Amap/Baidu Map）、音楽、電話、メッセージ、ポッドキャスト。Siri音声操作で完全ハンズフリー。
  \item \textbf{Android Auto}：Androidユーザー向けの対応策。機能はCarPlayに類似、Google Maps、Spotifyなどに対応。国内はGoogleサービス制限により一部機能使用不可。
  \item \textbf{Huawei HiCar}：Huawei自社開発方案、国内車種のサポート率が高い。ナビ·音楽のほか、Huaweiビデオ通話、駐車映像統合などの独自機能も使える。
  \item \textbf{Baidu CarLife}：クロスプラットフォーム方案、iOSとAndroidの両方で使用可能。機能は少ないが互換性が広く、多くの旧型車種はCarLifeのみ対応。
\end{itemize}

\begin{tcbwarn}{セキュリティに関する注意}
どのカースル連携方案でも、運転中はできるだけ音声操作を使うべきだ。ナビは事前に目的地を設定し、音楽は事前にプレイリストを選んでおき、走行中に車載画面を操作しない。

\end{tcbwarn}

\booksection{五、会議室キャスト実戦——もうバタバタしない}

会議はキャストの最も高頻度なシーンであり、最も問題が起きやすいシーンでもある。以下は実戦の経験則：

\begin{itemize}
  \item \textbf{事前到着してテスト}：会議開始直前にキャスト接続を試すのは絶対に避ける。10分前に到着し、接続方法をテストする。Wi-Fiキャストの失敗確率は想像より高い。
  \item \textbf{有線方案をPlan Bとして準備}：バッグにHDMI変換ケーブルを入れておく。ワイヤレスキャストが繋がらない時に有線が命綱になる。
  \item \textbf{通知プレビューをオフにする}：キャスト中、WeChatメッセージが会議室の大画面に表示される。キャスト前に設定で"ロック画面通知のプレビュー"をオフにするか、"プレゼンモード"をオンにする（一部スマホ対応、キャスト時に通知を自動非表示）。
  \item \textbf{"プレゼン文書"を使う而非"スマホ操作"}：PPTを展示する場合、PPTをスマホに保存してWPSで全画面再生する方が、スマホ画面全体をキャストするよりよい。よりプロフェッショナルでスムーズ。
\end{itemize}

\begin{tcbpractice}{実践：スマホからテレビへのキャストを一度完了する}
\begin{enumerate}
  \item テレビで「画面ミラーリング」または「AirPlay」機能を見つけてオン。
  \item スマホで写真または動画を一枚開く。
  \item ワイヤレスキャストを試す（iPhoneは画面ミラーリング、Androidはワイヤレスキャストボタン）。
  \item 失敗した場合、テレビが対応プロトコルをサポートしているか確認、またはDLNAプッシュを試す（アルバム/動画Appで"キャスト"ボタンをタップ）。
  \item キャスト成功後、スマホで他の操作をして、大画面が同期表示されるか確認する。
\end{enumerate}

\end{tcbpractice}

\begin{tcbsummary}{本章のまとめ}
\begin{itemize}
  \item ワイヤレスキャストの三大陣営：AirPlay（Apple）、Miracast（Android）、DLNA（クロスプラットフォーム·メディアプッシュ）
  \item 有線キャストはより安定·ゼロ遅延、会議室シーンにHDMI変換ケーブルが必須
  \item Samsung DeX/HuaweiデスクトップモードでスマホがライトオフィスPCに変身
  \item カースル連携：CarPlay/Android Auto/HiCar/CarLifeの四大方案
  \item 会議室キャストの鉄則：事前テスト、Plan B準備、通知プレビューオフ
\end{itemize}

\end{tcbsummary}
